\section{Two Cyclic Structures}
\label{sec:method}

\subsection{Tasks}
\label{sec:prelim}
\label{sec:problem}
At each step $t$ the agent receives an observation $o_t\in\R^n$, updates a
memory $h_t=f_\phi(o_t,h_{t-1})$, and assigns action probabilities
$\pi_\phi(a\mid h_t)$. The true position is one of $k$ positions on a known
cycle, and the memory may represent a single position or a probability over
several.

In the \textbf{mod-$k$} task, each of the $L$ observations encodes a binary event bit
$b_t\in\{0,1\}$. A single decision at the final step earns $+1$ if it equals the
running count $a^\star = \big(\sum_{t=1}^{L} b_t\big) \bmod k$ and $-1$ otherwise.
Parity is the case $k=2$. Models train on a bounded range of lengths and are
evaluated on longer ones. Two control tasks remove the cyclic structure: in the
\textbf{cue} task a single informative bit appears at $t=1$, and in the
\textbf{monotone} task the agent reports whether the count exceeds $L/2$. On
clean inputs, the programmed counter $c_t=(c_{t-1}+b_t)\bmod k$ is an exact
oracle, so the question is whether a policy can learn this counter from reward
alone. Sections~\ref{sec:transfer} and~\ref{sec:feedback} introduce the valve
and missing-event tasks.

\subsection{An input-conditioned phase rotation}
\label{sec:arch}
If the $k$ positions are marked around a circle, increasing the count by one
rotates a point by $2\pi/k$. We therefore keep a complex memory $z\in\C^d$ and,
at each step, rotate each coordinate by an angle chosen from the observation
before adding a learned input term:
\begin{equation}
z \;\leftarrow\; z \odot e^{\,i\,\theta(o_t)} \;+\; \mathrm{inj}(o_t),
\qquad \theta(o)=W_\theta\, o,
\label{eq:evolve}
\end{equation}
where $\odot$ denotes elementwise multiplication and $\theta(o)\in\R^d$ holds
the rotation angles. The rotation preserves each coordinate's magnitude, so $k$
advances of $2\pi/k$ return a coordinate exactly to its starting point; the
additive term $\mathrm{inj}(o)\in\C^d$ lets the observation change the
amplitudes as well. To read out an action, the model maps the state linearly to
one complex amplitude $\psi_a$ per action and normalizes the squared magnitudes.
We call this the Born readout:
\begin{equation}
\psi_a = \textstyle\sum_j w_{aj}\, z_j \in \C,
\qquad
\pi(a) = \frac{|\psi_a|^2}{\sum_b |\psi_b|^2}.
\label{eq:born}
\end{equation}
An all-zero score vector falls back to the uniform distribution, and sampling
floors probabilities at $10^{-8}$. When one coordinate stays fixed as a reference
and another rotates, the squared magnitude contains a term comparing their
phases, and this term peaks at the action aligned with the accumulated rotation.
Equivalently, the model is a real $SO(2)$ rotation with quadratic features.

\begin{proposition}[Exact and length-invariant representation;
cf.\ \citealp{grazzi2025negative,karuvally2025aussm}]
\label{prop:exact}
A parameter setting of Equations~\eqref{eq:evolve} and \eqref{eq:born} with width
$d=2$ makes the greedy action $\arg\max_a \pi(a)$ equal the true count $a^\star$ for
every length $L$ and every input sequence $b_1,\dots,b_L$.
\end{proposition}

The proof is in Appendix~\ref{app:proofs}. The proposition shows that exact
counting is representable; whether training finds such a setting is an empirical
question. A classical relative-phase score
$\operatorname{Re}(\overline{z_0}z_1\omega^{-a})$ selects the same action.
Appendix~\ref{app:proofs} also discusses the readout's global-phase invariance,
and Appendix~\ref{app:contractive} shows that uniformly contractive memories
cannot keep a positive decoding margin over long common suffixes.

\label{sec:stab}
\myparagraph{Training stabilizers and angular error.}
In the curriculum experiments we renormalize $z\leftarrow z/\lVert z\rVert_2$
after each step and add the penalty
$\lambda\langle\sin^2(\tfrac{k}{2}\theta(o_t))\rangle$ to pull angles toward
multiples of $2\pi/k$. Neither operation enforces exact angles, and a nonzero
injection can also shift the phase. For the two-channel construction without
injection, per-step phase errors of at most $\epsilon$ and an initial error $\epsilon_0$
give angular error at most $\eta=\epsilon_0+L\epsilon$ after $L$ steps. Nearest-residue
decoding stays correct for $\eta<\pi/k$, and with score perturbations of at most
$\delta$ it suffices that $\cos\eta-\cos(2\pi/k-\eta)>\delta$. Because the admissible error shrinks as $1/k$,
larger cycles are harder to track with a learned angle.

\subsection{A cyclic filter for uncertain increments}
\label{sec:filter}
A rotation carries a single count. Once an event may be missing or misread,
the model needs a probability for each position instead. Starting at zero, for
example, one missing turn leaves probability $0.50$ on position zero and $0.25$ on
each of positions one and three; a pointer reading of one then makes position one
more likely. A cyclic filter carries out these two steps as a prediction followed
by a correction:
\begin{align}
\tilde b_t &= \sum_{d} q_t(d)\, P_d\, b_{t-1},\label{eq:filter-predict}\\
b_t(s) &\propto \tilde b_t(s)\, \ell_t(y_t\mid s).\label{eq:filter-correct}
\end{align}
Here $P_d$ shifts the belief $b$ around the cycle by $d$ positions, a learned
event encoder supplies the probability $q_t(d)$ of each increment, and $\ell_t$
is the likelihood of the pointer reading $y_t$ at position $s$; without a reading,
$b_t=\tilde b_t$. Because the shifts are exact, only the event probabilities and
the pointer likelihood need to be learned. The categorical counter of
Section~\ref{sec:noisy} is the special case without pointer readings. For a
continuous increment $u_t$, the shift multiplies the belief's Fourier
coefficients by $e^{-inu_t}$, so the filter applies the phase rotation of
Section~\ref{sec:arch} to a whole distribution (Section~\ref{sec:imu}).
